% 第 1 章骨架：导引框已定，正文待写。写作计划见 docs/PLAN.md“第 1 章”，写法见 docs/STYLE.md。
\chapter{一个问题引出整门课}\label{ch:intro}

\begin{chapterguide}
\guidequestion{明朝宝钞、魏玛马克和津巴布韦元都崩溃了；日本政府债务超过 GDP 的两倍，却长期低息、低通胀；1998 年的俄罗斯和 2022 年的斯里兰卡违约了。差别在哪里？要回答它，需要知道什么？}
\begin{guidegoals}
\item 说出全书的主线：每一次理论革命都起于一场旧理论解释不了的危机，新理论又被下一场危机检验；
\item 用“四层读法”（会计恒等式、行为关系、均衡与预期的闭合、数据与政策观测量）拆解一个宏观论断；
\item 说明本币债与外币债、固定汇率与浮动汇率怎样改变一个国家“付得起”的含义；
\item 找到全书十一个部分各自回答的问题，并据此规划自己的阅读路线。
\end{guidegoals}
\guideprereq{无，从这里开始。}
\end{chapterguide}

\chapterstub
